\chapter{Metodologia}

Neste capítulo, serão apresentados os métodos utilizados para o desenvolvimento da análise multivariada, sendo estes: Limpeza de dados, Integração de dados, Seleção de dados, Transformação de dados, Mineração de dados, Avaliação de padrões e Apresentação do conhecimento.


Limpeza de dados:

1 - Dentre a definição das informações relevantes para a análise do câncer de mama, destaca-se dados públicos relacionados a fatores demográficos e socioeconômicos do estado MG, considerando apenas mulheres.
2 - Após a primeira restrição das áreas de interesse, analisamos as bases de dados disponíveis, identificando as variáveis presentes e confabilidade dos dados presentes. 
IFDM(Índice FIRJAN de Desenvolvimento Municipal) e Integrador RHC (Registro Hospitalar do Câncer) como fontes de dados para a análise.

Integração de dados:

3 - Com base nessa análise, foram definidas as bases de dados disponíveis do INCA (Instituto Nacional de Câncer), GOV (Governo do Estado de Minas Gerais), IPEA(Instituto de Pesquisa Econômica Aplicada)
Os dados podem ser correlacionados pelo código do município.

Seleção de dados:

4-A partir disso, partimos para a definição dos atributos de interesse, que foram definidos a partir de uma análise exploratória das bases de dados e filtros pré selecionados, buscando identificar 
as variáveis mais relevantes considerando aspectos mencionados na etapa anterior e que possibilitariam uma análise da série temporal desses conjuntos.
Sendo assim, defini-se os seguintes atributos de interesse:

    Atributos Demográficos:
    - Quantidade de hospitais habilitador em minas gerais, em 2025. 
    - Taxa de incidência de câncer de mama por município, analisandos por ano de diagnóstico, escolaridade, faixa etária e raça cor no período de 1985 até 2023
    - População por município de 0 até 75 anos ou mais, nos anos de 1970, 1980, 1991, 2000, 2010 e 2022 (anos censos)

    Atributos Socioeconômicos:

    - Taxa de mortalidade por câncer de mama por município, pelo período de 1979 até 2024
    - Índice de desenvolvimento humano municipal (IDHM) geral, de educação, saúde e emprego e renda, por município, de 2013 até 2024

Transformação de dados:

Inicialmente começou o processo de coleta dos dados, com o uso da linguagem Python, ver~sao 3.12, com técnicas de web scraper para extração e automatização da coleta. por meio de scripts simples
utilizando ferramentas como BeautifulSoup e Selenium, para acessar as páginas de interesse, extrair os dados relevantes e armazená-los em um formato estruturado para análise posterior, baixando arquivos em fomato csv

Aqui, os dados apresentados foram submetidos a uma série de tratamentos para ajuste e padronização, dentre os filtros aplicados temos:

    Normalização das nomeclaturas dos municípios e hospitais
    Padronização da série temporal para de 10 em 10 anos
    Filtragem das informações para apenas os municípios de Minas Gerais
    Valores em branco ou ausentes permancem assim
    Construção das seguintes tabelas

    Tabela 1: Quantidade de mortes por câncer de mama por município em mulheres, pelo período de 1979 até 2024, em que a linha representa a cidade e as colunas os anos, e o valor agregado bruto do número de mortes
    Tabela 2: Índice de desenvolvimento humano municipal (IDHM) geral, de educação, saúde e emprego e renda, por município, de 2013 até 2024, em que as linhas representam os municípios e as colunas os anos, e o valor agregado bruto do IDHM
   onde temos Desenvolvimento
        Crítico

        resultados
        inferiores a
        0,4 ponto

        
        Desenvolvimento
        Baixo

        resultados entre
        0,4 e 0,6 ponto

        
        Desenvolvimento
        Moderado

        resultados entre
        0,6 e 0,8 ponto

        
        Desenvolvimento
        Alto

        resultados
        superiores a
        0,8 ponto

 * A ausência de dados será identificada com um asterisco.

    tabela 3: Taxa de incidência de câncer de mama por município, em que a linha é o municíoo, e as colunas analisandos por ano de diagnóstico, escolaridade, faixa etária e raça cor no período de 1985 até 2023. Agregado a esta tabela, contém as informações sobre a população municipal, e os hospitais habilitados

Mineração de dados:



Avaliação de padrões: 

Dessa forma, encontrar o perfil epidemológico do câncer de mama em Minas Gerais, identificando as regiões mais vulneráveis, os fatores de risco associados e as possíveis lacunas no acesso ao diagnóstico e tratamento. A partir disso, será possível subsidiar políticas públicas e estratégias de prevenção direcionadas às populações mais vulneráveis, visando a redução da incidência e mortalidade pela doença.

Apresentação do conhecimento:

A metodologia adotada segue os princípios adotados pelo KDD.


Falta definir a parte de algoritmos e métricas de avaliação

% \section{Introduction}
% In this example, several keywords\index{keywords} will be used which
% are important and deserve to appear in the Index\index{Index}. Terms like generate\index{generate} and some\index{others} will also
% show up. Terms in the index can also be nested \index{Index!nested}

% \section{Second section}
% This second section\index{section} may include some special word,
% and expand the ones already used\index{keywords!used}.